\section{Design}\label{sec:design}

\sys{} is a single-queue gRPC broker (\texttt{Submit}, \texttt{Pull}, \texttt{Ack}, \texttt{Nack}) with three
priority levels, retries, a dead-letter queue and at-least-once delivery while the process runs. Consumers are
tier-blind: the broker decides which message each \texttt{Pull} returns. State is in memory.

\paragraph{Keys and prediction.}
A message's key is the producer id, optionally joined to a producer-supplied \texttt{job\_type} header. Each key has a
decaying histogram of handler durations; there is also a global histogram. On \texttt{Ack}, the consumer's reported
\texttt{processing\_time\_ms} is added to the key's histogram; only successful, non-replayed durations are learned.
At ingress, \texttt{route\_message()} looks up the key's median and maps it to a level using boundaries derived from
global quantiles. A key with too few samples, a stale history, or high spread takes the default (middle) level.
Per-key state is bounded. A lookup takes about 2\,$\mu$s.

\paragraph{Aging.}
Strict priority starves long jobs, so the broker promotes messages that wait past a threshold (default 5\,s). This
bounds long-job waits, but under backlogs lasting minutes or hours it promotes almost everything to the top level
and erases the predictor's ordering (Section~\ref{sec:results}). For that regime we added an opt-in design: worker
time is shared across levels by weights (e.g.\ 8:3:1), and promotion into a level pauses while that level is behind
its share (\emph{pausing aging}).

\paragraph{Modes.}
\emph{Static} routing uses a fixed default level and has no predictor. \emph{Shadow} mode records predictions but
routes statically, which isolates the cost of prediction. \emph{Predictive} mode routes by prediction. Default routing
remains static.
